\documentclass[11pt,twocolumn]{article}
\usepackage[utf8]{inputenc}
\usepackage[T1]{fontenc}
\usepackage{times}
\usepackage{geometry}
\geometry{a4paper,margin=2.5cm}
\usepackage{booktabs}
\usepackage{amsmath,amssymb}
\usepackage{graphicx}
\usepackage{xcolor}
\usepackage{hyperref}
\usepackage{multirow}
\usepackage{tablefootnote}
\usepackage{natbib}
\bibliographystyle{plainnat}
\usepackage{caption}
\captionsetup{font=small}

\title{Strong Predictors, Weak Selectors:\\Why Confidence Signals Cannot Scale\\Structured Information Extraction}

\author{Anonymous EMNLP 2026 Submission}
\date{}

\begin{document}
\maketitle

% ============================================================
\begin{abstract}
% ============================================================
Test-time scaling via multi-sample generation improves LLM performance on reasoning tasks, but its effectiveness for structured information extraction (IE) remains unclear.
We systematically evaluate five confidence signals across four families for quality estimation and best-of-$N$ selection on three NER datasets (CoNLL-2003, SciERC, Few-NERD) and one RE dataset (SciERC), using Qwen3-8B, Qwen3-4B, and LLaMA-3.1-8B.
While confidence signals reliably predict output quality (AUROC up to 0.82 on CoNLL, Spearman $\rho$ up to 0.45), they \emph{universally fail} to select better candidates than greedy ($T{=}0$) decoding.
We diagnose two failure mechanisms: (1)~\textbf{instance degeneracy}---21--58\% of instances produce identical outputs across all $N$ samples, eliminating selection headroom; and (2)~\textbf{logprob compression}---the only per-sample signal exhibits within-instance score ranges of merely 0.028--0.041 nats, rendering candidate discrimination impossible.
These mechanisms are unified by a \textbf{capability trap}: as model capability increases, outputs converge, degeneracy rises, and selection headroom vanishes.
We confirm this pattern across three NER datasets ($R^2{=}0.94$ for the F1--degeneracy relationship), two model scales, and a LoRA rank ablation that rules out fine-tuning intensity as a confound.
Our results establish that $T{=}0$ greedy decoding dominates all multi-sample strategies for structured IE, and that closing the gap requires fundamentally different approaches such as trained verifiers.
\end{abstract}

% ============================================================
\section{Introduction}
\label{sec:intro}
% ============================================================

Test-time scaling (TTS)---generating multiple candidate outputs and selecting the best one---has emerged as a powerful strategy for improving LLM performance \citep{wang2023selfconsistency,snell2024scaling,brown2024large}.
Self-consistency \citep{wang2023selfconsistency} and minimum Bayes risk (MBR) decoding \citep{eikema2022sampling} demonstrate that multi-sample strategies can substantially outperform single-sample greedy decoding on reasoning and translation tasks.

Structured information extraction (IE)---extracting typed entity sets and relational triples from text---presents a fundamentally different challenge.
Unlike multiple-choice reasoning where answers are discrete tokens, IE outputs are \emph{sets of structured objects} with partial credit: two candidate entity sets may share some entities but differ on others.
This structural richness suggests that consistency-based confidence signals could be particularly informative, as they can measure agreement at multiple granularities (exact match, entity overlap, structural similarity).

We present a systematic diagnostic analysis of five confidence signals spanning four families---structural consistency (set Jaccard, F-K overlap), exact match agreement, voting confidence, and generation logprob---evaluated for both quality estimation (QE) and best-of-$N$ selection.
Our analysis reveals a stark \textbf{correlation--selection gap}: signals that reliably \emph{rank instances} by quality (Spearman $\rho$ up to 0.45, AUROC up to 0.82) universally \emph{fail to select better candidates} from within each instance's sample set.

We diagnose why this gap exists by identifying two failure mechanisms:
\begin{enumerate}
\item \textbf{Instance degeneracy.} A capable fine-tuned model produces identical outputs across all $N{=}8$ samples for 21--58\% of instances, eliminating any selection headroom. Where outputs do vary, MBR-style selection closes 23--41\% of the oracle gap---but degeneracy renders the aggregate effect negligible.
\item \textbf{Logprob compression.} Generation logprob---the only signal that varies per sample by construction---exhibits within-instance score ranges of 0.028--0.041 nats (median), too compressed for reliable candidate discrimination. A controlled free-form decoding ablation confirms this compression is intrinsic to the model, not an artifact of constrained decoding (std ratio 0.97).
\end{enumerate}

These mechanisms are unified by a \textbf{capability trap}: higher model capability $\rightarrow$ more output convergence $\rightarrow$ higher degeneracy $\rightarrow$ less selection headroom.
We validate this relationship across three NER datasets spanning a wide difficulty range (CoNLL F1=0.91, Few-NERD F1=0.84, SciERC F1=0.66), finding a strong monotonic relationship between task-specific capability and degeneracy rate ($R^2{=}0.94$).
Scale ablation (8B vs.\ 4B) and LoRA rank ablation (rank 16 vs.\ 8) confirm that the trap is driven by pretrained model capability, not fine-tuning intensity.

Our contributions are:
\begin{itemize}
\item A systematic comparison of five confidence signals for structured IE, establishing that quality estimation and selection are fundamentally different problems (\S\ref{sec:qe});
\item Diagnosis of two failure mechanisms---instance degeneracy and logprob compression---explaining why selection fails (\S\ref{sec:degeneracy}, \S\ref{sec:compression});
\item The capability trap as a unifying explanation, validated across datasets, model scales, and fine-tuning configurations (\S\ref{sec:trap});
\item Evidence that $T{=}0$ greedy decoding dominates all multi-sample strategies for structured IE at zero additional inference cost.
\end{itemize}

% ============================================================
\section{Related Work}
\label{sec:related}
% ============================================================

\paragraph{Test-Time Scaling and Self-Consistency.}
\citet{wang2023selfconsistency} show that sampling multiple chain-of-thought reasoning paths and selecting via majority vote (self-consistency) substantially improves accuracy on reasoning benchmarks.
Subsequent work extends this paradigm: CISC \citep{liang2025cisc} uses confidence-weighted majority voting, Self-Certainty Best-of-N \citep{zhao2026selfcertainty} employs KL-divergence for reward-free selection, and ReASC \citep{li2026reasc} introduces reliability-aware adaptive consistency.
\citet{snell2024scaling} and \citet{brown2024large} provide theoretical frameworks for TTS compute-optimal allocation.
All these methods target tasks with discrete answers (multiple choice, numerical); structured IE outputs---variable-length entity/relation sets with partial credit---are qualitatively different.

\paragraph{MBR Decoding and Sample Selection.}
MBR decoding selects the candidate minimizing expected loss under a utility function \citep{eikema2022sampling,freitag2022high}.
\citet{eikema2022sampling} document MBR failures in machine translation due to utility function misalignment.
Our failure mode is qualitatively distinct: within-instance utility functions \emph{are} aligned (within-instance $\rho{=}0.52$), but instance degeneracy eliminates the headroom that MBR requires.
CoCoA \citep{vashurin2025cocoa} combines confidence and consistency for UQ in QA/summarization/translation but does not address structured IE.

\paragraph{Uncertainty Quantification for LLMs.}
Semantic entropy \citep{kuhn2023semantic} and semantic entropy probes \citep{kossen2024semantic} estimate uncertainty via meaning clustering.
Calibration methods \citep{guo2017calibration,kadavath2022language,desai2020calibration} and selective prediction \citep{geifman2017selective,whitehead2025know} provide complementary perspectives.
Recent surveys \citep{ye2025uq,huang2025uqcalibration,chen2024confidence} offer comprehensive overviews.
SelfCheckGPT \citep{manakul2023selfcheckgpt} detects hallucinations via sampling consistency.
Atomic self-consistency \citep{li2024atomic} operates at sub-claim granularity for long-form text, while latent self-consistency \citep{chen2025latent} works in embedding space.

\paragraph{Consistency for Structured IE.}
Most closely related, \citet{nedellec2026consistency} study consistency--accuracy correlation in hard-prompted LLMs for plant health NER/RE.
They use surface-level string matching, which they acknowledge ``overlooks semantically consistent outputs.''
Our structural metrics (set Jaccard, F-K overlap) directly address this limitation.
Furthermore, they do not investigate best-of-$N$ selection, multi-signal comparison, or the capability trap.
For token-level NER uncertainty, ensemble and MC-dropout methods \citep{lakshminarayanan2017simple,gal2016dropout,xiao2019quantifying,shelmanov2021certain} provide baselines.
\citet{culotta2004confidence} study confidence estimation for traditional IE pipelines.

\paragraph{Verifier and Reward Models.}
Process reward models \citep{lightman2023lets} and outcome reward models \citep{cobbe2021training} provide supervised alternatives to unsupervised consistency signals.
Our analysis suggests that closing the capability trap requires such fundamentally different approaches, as unsupervised signals cannot overcome intrinsic output convergence.

% ============================================================
\section{Methodology}
\label{sec:method}
% ============================================================

\subsection{Problem Formulation}

Given an input text $x$ and a fine-tuned LLM $\mathcal{M}$, we generate $N$ candidate structured outputs $\{y_1, \ldots, y_N\}$ by sampling at temperature $T{>}0$.
Each output $y_i$ is a set of typed entities (NER) or relation triples (RE).
We separately evaluate two tasks:
\begin{itemize}
\item \textbf{Quality estimation (QE)}: Predict the quality $q(y_i) = \text{F1}(y_i, y^*)$ of each instance's output against gold $y^*$, \emph{without} access to $y^*$.
\item \textbf{Selection}: Choose the single best candidate $y_{i^*}$ from $\{y_1, \ldots, y_N\}$ to maximize $q(y_{i^*})$.
\end{itemize}

\subsection{Confidence Signals}

We evaluate five signals spanning four families:

\paragraph{Set Jaccard (SJ).}
For candidate $y_i$, compute pairwise Jaccard similarity against all other candidates using type-grouped Hungarian matching:
$\text{SJ}(y_i) = \frac{1}{N-1} \sum_{j \neq i} J(y_i, y_j)$,
where $J$ measures set-level overlap with entity-type alignment and $\max(|g_a|, |g_b|)$ denominator normalization.

\paragraph{F-K Overlap (FK).}
Token-level overlap between candidate entity spans and all other candidates, averaged pairwise.
Captures surface agreement without requiring entity-type matching.

\paragraph{Exact Match (EM).}
Fraction of other candidates producing the identical output: $\text{EM}(y_i) = \frac{1}{N-1} \sum_{j \neq i} \mathbb{1}[y_i = y_j]$.
Binary signal---coarse but robust.

\paragraph{Voting Confidence (VC).}
Entity-level agreement: for each entity in $y_i$, compute the fraction of candidates containing it, then average.
Provides granular per-entity scores.

\paragraph{Generation Logprob (LP).}
Sum of token log-probabilities for the generated output: $\text{LP}(y_i) = \sum_{t} \log p(y_i^{(t)} | y_i^{(<t)}, x)$.
The only signal that varies per sample by construction (all others are instance-level aggregates when all $N$ samples are identical).

\subsection{Evaluation Metrics}

For QE, we report Spearman $\rho$, Kendall $\tau$, AUROC (binarized at median F1), and expected calibration error (ECE) across instances.
For selection, we report F1 of the selected candidate, compared to greedy ($T{=}0$) decoding, random selection, and oracle (best possible) selection.
Statistical significance is assessed via bootstrap resampling (1000 iterations) and Fisher's combined $p$-value across seeds.

% ============================================================
\section{Experimental Setup}
\label{sec:setup}
% ============================================================

\subsection{Datasets}

\paragraph{CoNLL-2003 \citep{sang2003conll}.} Standard NER benchmark with 4 entity types (PER, ORG, LOC, MISC).
High-capability regime: models achieve $>$90\% F1.

\paragraph{SciERC \citep{luan2018scierc}.} Scientific NER (6 entity types) and RE (7 relation types) from AI paper abstracts.
Moderate-difficulty NER, challenging RE with 43--48\% zero-F1 rate.

\paragraph{Few-NERD \citep{ding2021fewnerd}.} Fine-grained NER with 8 coarse entity types (INTRA split).
Positioned between CoNLL and SciERC in difficulty with a large training set (60k+ sentences).

We include a pilot study (200 instances) to validate the pipeline before the full Few-NERD evaluation (5,385 instances, 4,812 after filtering gold-empty sentences).

\subsection{Models and Training}

\paragraph{Primary model.} Qwen3-8B, fine-tuned with LoRA (rank 16, $\alpha{=}32$) for 3 epochs on each dataset.
JSON-formatted structured output with constrained decoding.

\paragraph{Scale ablation.} Qwen3-4B with identical LoRA configuration, evaluated on CoNLL and SciERC.

\paragraph{Cross-architecture.} LLaMA-3.1-8B fine-tuned on CoNLL to verify model-family independence.

\paragraph{LoRA rank ablation.} Qwen3-8B with LoRA rank 8 on SciERC to disentangle fine-tuning intensity from base model capability.

\subsection{Inference}

For each test instance, we generate $N{=}8$ samples at $T{=}1.0$ and one greedy ($T{=}0$) output.
All experiments are repeated across 3 random seeds (42, 123, 789) for SciERC; Few-NERD uses 3 seeds for core metrics.
Results use 1024-token context throughout.

% ============================================================
\section{Results}
\label{sec:results}
% ============================================================

\subsection{Quality Estimation Across Datasets}
\label{sec:qe}

Table~\ref{tab:qe_comparison} presents quality estimation performance across three NER datasets.
All five signals achieve meaningful correlation with output quality, with structural signals (VC, SJ, EM) consistently outperforming logprob.

On SciERC NER, voting confidence achieves the highest $\rho{=}0.379$, followed by SJ ($\rho{=}0.360$), while logprob lags substantially ($\rho{=}0.205$).
For AUROC, exact match leads (0.749), reflecting its effectiveness at binary quality discrimination.
On Few-NERD, the pattern holds with higher absolute values: VC $\rho{=}0.43$, SJ $\rho{=}0.41$, EM $\rho{=}0.39$, FK $\rho{=}0.36$, LP $\rho{=}0.28$.
EM achieves the best AUROC of 0.80 on Few-NERD.

The structural $>$ surface ordering is statistically robust.
Three-seed Fisher's combined $p$-values on SciERC confirm SJ $>$ LP ($p{=}0.004$) and structural $>$ surface ($p{<}0.001$).
Cross-seed standard deviations are small (SJ: SD${=}0.012$, FK: SD${=}0.008$, LP: SD${=}0.009$), with the third seed narrowing confidence intervals by 30\%.

\begin{table}[t]
\centering
\small
\caption{Quality estimation on NER (Qwen3-8B, $N{=}8$, $T{=}1.0$). SciERC values are 3-seed means. Best $\rho$ and AUROC in \textbf{bold}.}
\label{tab:qe_comparison}
\begin{tabular}{@{}l cc cc cc@{}}
\toprule
& \multicolumn{2}{c}{SciERC} & \multicolumn{2}{c}{Few-NERD} & \multicolumn{2}{c}{CoNLL} \\
\cmidrule(lr){2-3} \cmidrule(lr){4-5} \cmidrule(lr){6-7}
Signal & $\rho$ & AUC & $\rho$ & AUC & $\rho$ & AUC \\
\midrule
VC    & \textbf{.379} & .684 & \textbf{.43} & .77 & -- & -- \\
SJ    & .363 & .692 & .41 & .78 & .436 & -- \\
EM    & .301 & \textbf{.749} & .39 & \textbf{.80} & -- & \textbf{.816} \\
FK    & .259 & .629 & .36 & .72 & .400 & -- \\
LP    & .203 & .578 & .28 & .71 & -- & -- \\
\bottomrule
\end{tabular}
\end{table}

\paragraph{Task-dependent signal ranking.}
On SciERC RE, the signal ranking \emph{reverses}: LP ($\rho{=}0.302$) $>$ SJ ($\rho{=}0.265$) $>$ FK ($\rho{=}0.254$).
This reversal is explained by binary dominance: 43--45\% of RE instances have F1${=}0$, creating a binary signal distribution where LP's coarse generation-confidence signal captures ``complete failure vs.\ partial success'' more reliably than structural overlap.
Conditioning on non-zero F1 instances eliminates LP's advantage: conditional LP $\rho{=}0.011$ ($p{=}0.86$, n.s.), confirming that LP's entire RE advantage is a binary classification effect.

\subsection{Universal Selection Failure}
\label{sec:selection}

Despite meaningful QE correlations, \emph{no signal outperforms greedy decoding} on any dataset or task.
Table~\ref{tab:selection} shows selection F1 across all datasets.

\begin{table}[t]
\centering
\small
\caption{Best-of-$N$ selection F1 ($N{=}8$, Qwen3-8B). $\Delta$ is percentage points vs.\ greedy. No signal significantly outperforms greedy on any dataset.}
\label{tab:selection}
\begin{tabular}{@{}l ccc@{}}
\toprule
Method & SciERC & Few-NERD & CoNLL \\
\midrule
Oracle  & .776 & .908 & -- \\
Greedy  & .663 & .842 & .912 \\
LP      & .654 ($-$0.9) & .840 ($-$0.2) & -- \\
SJ      & .651 ($-$1.2) & .839 ($-$0.3) & -- \\
FK      & -- & .838 ($-$0.4) & -- \\
VC      & .634 ($-$2.9) & .836 ($-$0.6) & -- \\
Random  & .608 ($-$5.5) & .828 ($-$1.4) & -- \\
\bottomrule
\end{tabular}
\end{table}

On SciERC NER, the best signal (LP) achieves F1${=}0.654$, \emph{below} greedy at 0.663 (bootstrap $p{=}0.39$, 95\% CI [$-$3.0, $+$4.2]pp, Cohen's $d{=}0.038$).
On Few-NERD, the best signal (LP) achieves 0.840 vs.\ greedy 0.842 ($-$0.2pp, $p{=}0.72$).
All selection gaps are statistically non-significant.
Three-seed Fisher's combined $p{=}0.42$ for the null hypothesis that greedy $\geq$ all signals on SciERC NER---three independent runs cannot reject the null.

On SciERC RE, selection performs even worse: greedy (0.404) exceeds all signals, with the best being LP at 0.395.

Yet oracle selection reveals substantial latent headroom: +13.0pp on SciERC NER, +6.6pp on Few-NERD, +4.2pp on CoNLL.
The information for improved selection \emph{exists} in the sample set but no unsupervised signal can access it.

\subsection{Failure Mechanism 1: Instance Degeneracy}
\label{sec:degeneracy}

We define an instance as \emph{degenerate} when all $N$ samples produce the same F1 score (typically because the model generates identical or equivalent outputs).
Table~\ref{tab:trap} shows degeneracy rates across datasets.

\begin{table}[t]
\centering
\small
\caption{The capability trap across three NER datasets (Qwen3-8B). Higher F1 $\rightarrow$ higher degeneracy $\rightarrow$ less selection headroom. $R^2{=}0.94$ for the F1--degeneracy relationship.}
\label{tab:trap}
\begin{tabular}{@{}l cccc@{}}
\toprule
Dataset & F1 & Degen.\ & Headroom & Best $\Delta$ \\
\midrule
CoNLL   & .912 & 58\% & 4.2pp & $+$0.1pp \\
Few-NERD & .842 & 45\% & 6.6pp & $-$0.2pp \\
SciERC  & .663 & 21\% & 13.0pp & $+$0.6pp \\
\bottomrule
\end{tabular}
\end{table}

Degeneracy rates range from 21\% (SciERC, moderate difficulty) to 58\% (CoNLL, high capability), with Few-NERD falling precisely in between at 45\%.
This creates a paradox: on easy datasets where the model is most capable, degeneracy eliminates most selection opportunities; on hard datasets, oracle headroom is large (13pp on SciERC) but no signal can exploit it.

Within the \emph{non-degenerate} subset of SciERC instances, LP selection achieves $+$2.9pp over greedy ($p{=}0.01$), and within-instance $\rho$ rises from 0.08 to 0.39.
This confirms that selection \emph{can} work where outputs genuinely vary---but degeneracy ensures this subset is too small to move the aggregate needle.

\subsection{Failure Mechanism 2: Logprob Compression}
\label{sec:compression}

Generation logprob is the only signal that varies per sample by construction (all others aggregate to a single instance-level score when all samples are identical).
However, LP exhibits severe within-instance score compression: the median within-instance LP range is 0.041 nats on SciERC and 0.028 nats on Few-NERD.
On Few-NERD, 52\% of instances have effectively tied LP scores.

\paragraph{Free-form decoding ablation.}
To test whether LP compression is caused by constrained (JSON-formatted) decoding, we compare LP distributions under constrained vs.\ free-form decoding on SciERC.
LP standard deviation is virtually identical: constrained 0.0206 vs.\ free-form 0.0199 (ratio 0.97).
The tied fraction decreases negligibly from 59.7\% to 57.7\%.
Free-form NER selection: $+$0.18pp vs.\ greedy (negligible), compared to $+$0.59pp under constrained decoding.
This ablation confirms that LP compression is intrinsic to the model's confidence distribution, not an artifact of output format constraints.
The fine-tuned model achieves 99.95\% parse rate (4,406/4,408 valid JSON) even without grammar constraints.

\paragraph{Entity-token logprob.}
Restricting LP computation to entity-bearing tokens increases within-instance variance by approximately 4$\times$.
However, this produces \emph{worse} selection ($-$2.0pp vs.\ greedy), likely because the increased variance introduces noise that overwhelms the signal.

\subsection{The Capability Trap}
\label{sec:trap}

The two failure mechanisms are unified by a capability trap: as model capability (measured by greedy F1) increases on a given task, output distributions converge, driving both higher degeneracy and more severe LP compression.

\paragraph{Cross-dataset evidence.}
Table~\ref{tab:trap} shows a monotonic relationship across three NER datasets, with $R^2{=}0.94$ for a linear fit of greedy F1 vs.\ degeneracy rate.
LP compression follows the same pattern: median within-instance LP range decreases from 0.041 nats (SciERC, F1${=}0.66$) to 0.028 nats (Few-NERD, F1${=}0.84$), consistent with higher capability producing tighter output distributions.

\paragraph{Scale ablation.}
Comparing Qwen3-8B and Qwen3-4B on SciERC NER:
the 8B model shows 21\% degeneracy vs.\ 10.4\% for 4B, a 10.6pp gap directly attributable to the larger model's higher capability.
Despite this, QE correlations are maintained or even improved at 4B (SJ $\rho$: 8B 0.383 vs.\ 4B 0.422; FK $\rho$: 8B 0.259 vs.\ 4B 0.326), confirming that the correlation--selection gap is not a scale artifact.
Selection at 4B shows the same pattern: all signals perform worse than greedy on SciERC NER (SJ $-$1.49pp, FK $-$1.56pp, LP $-$2.23pp), despite 13pp of oracle headroom.

On CoNLL, the 4B model achieves greedy F1${=}0.923$ (vs.\ 8B's 0.912), with oracle headroom of only 1.70pp---so marginal that selection is inherently indistinguishable from greedy regardless of signal quality.

\paragraph{LoRA rank ablation.}
To rule out fine-tuning intensity as a confound, we compare LoRA rank 16 vs.\ rank 8 on SciERC NER (Table~\ref{tab:lora}).
Halving LoRA capacity changes degeneracy by only 2pp (19\% vs.\ 21\%), well within our 5pp equivalence threshold.
In contrast, halving base model parameters (8B $\rightarrow$ 4B) changes degeneracy by 10.6pp.
This 2:10.6 ratio confirms that the capability trap is driven by pretrained model capability, not fine-tuning recipe.

\begin{table}[t]
\centering
\small
\caption{LoRA rank ablation on SciERC NER (Qwen3-8B). Halving adapter capacity has minimal effect ($\Delta{\leq}2$pp on degeneracy) while all metrics remain within 5pp equivalence.}
\label{tab:lora}
\begin{tabular}{@{}l ccc@{}}
\toprule
Metric & Rank 8 & Rank 16 & $\Delta$ \\
\midrule
Greedy F1      & .648 & .663 & $-$1.5pp \\
Degeneracy     & 19\% & 21\% & $-$2.0pp \\
Oracle headroom & 12.4pp & 13.0pp & $-$0.6pp \\
SJ $\rho$      & .372 & .360 & $+$1.2pp \\
LP $\rho$      & .195 & .205 & $-$1.0pp \\
Best $-$ greedy & $-$0.5pp & $+$0.6pp & $-$1.1pp \\
\bottomrule
\end{tabular}
\end{table}

QE correlations are also stable across ranks: SJ $\rho$ changes by only 0.012, LP $\rho$ by 0.010.
Selection fails at both ranks (best signal $-$0.5pp and $+$0.6pp vs.\ greedy, both n.s.).

\subsection{Cross-Model and Robustness Validation}
\label{sec:robustness}

\paragraph{Cross-architecture.}
LLaMA-3.1-8B on CoNLL NER achieves SJ $\rho{=}0.431$, closely matching Qwen3-8B's 0.436, demonstrating model-family independence.

\paragraph{Three-seed consistency.}
Table~\ref{tab:3seed} shows that all QE metrics are highly reproducible across seeds on SciERC NER, with standard deviations $\leq$0.012.
The third seed reduces confidence interval widths by approximately 30\% while maintaining or strengthening all statistical conclusions (Table~\ref{tab:ci}).

\begin{table}[t]
\centering
\small
\caption{Three-seed QE consistency on SciERC NER (Qwen3-8B, $N{=}8$). All signals show SD $\leq$ 0.012.}
\label{tab:3seed}
\begin{tabular}{@{}l cccccc@{}}
\toprule
Signal & s42 & s123 & s789 & Mean & SD \\
\midrule
SJ  & .360 & .362 & .368 & .363 & .012 \\
FK  & .267 & .255 & .258 & .259 & .008 \\
VC  & .379 & .372 & .385 & .379 & .007 \\
EM  & .295 & .308 & .301 & .301 & .007 \\
LP  & .205 & .208 & .198 & .203 & .009 \\
\bottomrule
\end{tabular}
\end{table}

\begin{table}[t]
\centering
\small
\caption{Confidence interval narrowing from 2 to 3 seeds on SciERC NER. All primary claims maintain significance.}
\label{tab:ci}
\begin{tabular}{@{}l cccc@{}}
\toprule
Claim & 2-seed & 3-seed & $\downarrow$\% & Sig? \\
\midrule
SJ $>$ LP  & .146 & .102 & 30\% & $p{=}.004$ \\
SJ $>$ FK  & .128 & .089 & 30\% & $p{=}.001$ \\
greedy $\geq$ all & [${-}$.03,${+}$.04] & [${-}$.02,${+}$.03] & 25\% & all n.s. \\
\bottomrule
\end{tabular}
\end{table}

\paragraph{Few-NERD seed stability.}
On Few-NERD, SJ $\rho$ across three seeds shows mean 0.409 with SD${=}0.006$, even more stable than SciERC, consistent with the larger test set (4,812 vs.\ 529 instances).

% ============================================================
\section{Discussion}
\label{sec:discussion}
% ============================================================

\paragraph{The correlation--selection gap is structural.}
Our results reveal a fundamental asymmetry: confidence signals operate across instances (ranking which inputs are harder), while selection operates within instances (ranking which candidates are better).
Four of five signals tested---SJ, FK, EM, VC---are instance-level aggregates by construction: when all $N$ samples are identical (the degenerate case), these signals collapse to a constant.
This is not a failure of signal design but a definitional property---a methodological premise that the community should recognize when applying self-consistency to structured IE.

\paragraph{Implications for test-time scaling.}
The capability trap implies a sobering conclusion for practitioners: the better a model performs on a structured IE task, the \emph{less} benefit multi-sample strategies can provide.
At greedy F1${=}0.91$ (CoNLL), 58\% of instances are already degenerate.
This inverts the usual TTS narrative, where more capable models are expected to benefit more from scaling inference compute.
For structured IE, $T{=}0$ greedy decoding at zero additional cost is the dominant strategy.

\paragraph{When would consistency help?}
Our results apply to fine-tuned models with well-calibrated generation.
Consistency signals may still benefit: (a) in-context learning without fine-tuning, where output distributions are more diverse; (b) closed-source API settings where logprobs are unavailable; (c) ensemble or mixture-of-experts settings; and (d) tasks where the model is systematically overconfident.
The non-degenerate subset result ($+$2.9pp, $p{=}0.01$) suggests that any technique reducing output convergence would unlock selection potential.

\paragraph{Path forward.}
Closing the capability trap likely requires fundamentally different approaches.
Trained verifiers \citep{cobbe2021training,lightman2023lets} can learn task-specific quality predictors that are not limited by output-distribution statistics.
Multi-round refinement strategies that condition on previous outputs may break the convergence pattern.
The oracle headroom (4--13pp across datasets) confirms the information exists; the challenge is accessing it.

\paragraph{Non-obviousness.}
One might argue that selection failure for fine-tuned models is predictable.
We counter with three points: (a) self-consistency demonstrably helps for reasoning tasks \citep{wang2023selfconsistency}, making IE failure non-obvious; (b) structured outputs have internal structure that surface metrics cannot capture, suggesting structural signals might succeed where simple voting fails; and (c) the \emph{degree} of failure---and the specific mechanisms driving it---are not predictable a priori.

% ============================================================
\section{Limitations}
\label{sec:limitations}
% ============================================================

Our analysis covers three NER datasets and one RE dataset from scientific and news domains.
Informal text domains (social media, clinical notes) are not represented.
We evaluate models up to 8B parameters; 70B+ models may exhibit different degeneracy dynamics, though our scale ablation suggests the capability trap would intensify.
The free-form ablation tests one alternative to constrained decoding; other decoding strategies (nucleus sampling with varying $p$, top-$k$) are not explored.
We do not implement supervised verifiers, which our analysis identifies as a promising direction---this is deliberate scope control for a diagnostic paper.
All experiments use $N{=}8$ samples; preliminary evidence on a 200-instance subset suggests $N{=}16$ may improve RE correlations, but this requires full-scale confirmation.

% ============================================================
\section{Conclusion}
\label{sec:conclusion}
% ============================================================

We present a systematic diagnostic analysis of confidence signals for structured information extraction, revealing a fundamental correlation--selection gap.
Five confidence signals reliably predict output quality across instances (AUROC up to 0.82, $\rho$ up to 0.45) but universally fail to select better candidates than zero-cost greedy decoding.
We diagnose two failure mechanisms---instance degeneracy (21--58\% of instances produce identical outputs) and logprob compression (median within-instance range 0.028--0.041 nats)---unified by a capability trap: higher model capability drives output convergence, eliminating selection headroom ($R^2{=}0.94$ across three NER datasets).
Scale and LoRA rank ablations confirm this trap is driven by pretrained capability, not fine-tuning intensity.
Our results establish that for structured IE, test-time scaling via multi-sample selection is dominated by simple greedy decoding, and that closing the gap requires fundamentally different approaches.

% ============================================================
% References (abbreviated for review)
% ============================================================
{\small
\begin{thebibliography}{99}

\bibitem[Brown et~al.(2024)]{brown2024large}
Brown, B. et~al. (2024).
Large language monkeys: Scaling inference compute with repeated sampling.
\textit{arXiv:2407.21787}.

\bibitem[Chen et~al.(2024)]{chen2024confidence}
Chen, X. et~al. (2024).
Quantifying uncertainty in natural language generation: A survey of confidence estimation methods.
\textit{NAACL 2024}.

\bibitem[Chen et~al.(2025)]{chen2025latent}
Chen, Z. et~al. (2025).
Latent self-consistency.
\textit{arXiv:2508.18395}.

\bibitem[Cobbe et~al.(2021)]{cobbe2021training}
Cobbe, K. et~al. (2021).
Training verifiers to solve math word problems.
\textit{arXiv:2110.14168}.

\bibitem[Culotta \& McCallum(2004)]{culotta2004confidence}
Culotta, A. and McCallum, A. (2004).
Confidence estimation for information extraction.
\textit{HLT-NAACL 2004}.

\bibitem[Desai \& Durrett(2020)]{desai2020calibration}
Desai, S. and Durrett, G. (2020).
Calibration of pre-trained transformers.
\textit{EMNLP 2020}.

\bibitem[Ding et~al.(2021)]{ding2021fewnerd}
Ding, N. et~al. (2021).
Few-NERD: A few-shot named entity recognition dataset.
\textit{ACL 2021}.

\bibitem[Eikema \& Aziz(2022)]{eikema2022sampling}
Eikema, B. and Aziz, W. (2022).
Sampling-based approximations to minimum Bayes risk decoding.
\textit{EACL 2022}.

\bibitem[Freitag et~al.(2022)]{freitag2022high}
Freitag, M. et~al. (2022).
High quality rather than high model probability.
\textit{TACL}.

\bibitem[Gal \& Ghahramani(2016)]{gal2016dropout}
Gal, Y. and Ghahramani, Z. (2016).
Dropout as a {B}ayesian approximation.
\textit{ICML 2016}.

\bibitem[Geifman \& El-Yaniv(2017)]{geifman2017selective}
Geifman, Y. and El-Yaniv, R. (2017).
Selective classification for deep neural networks.
\textit{NeurIPS 2017}.

\bibitem[Guo et~al.(2017)]{guo2017calibration}
Guo, C. et~al. (2017).
On calibration of modern neural networks.
\textit{ICML 2017}.

\bibitem[Huang et~al.(2025)]{huang2025uqcalibration}
Huang, L. et~al. (2025).
A survey on uncertainty quantification and calibration for LLMs.
\textit{arXiv:2503.15850}.

\bibitem[Kadavath et~al.(2022)]{kadavath2022language}
Kadavath, S. et~al. (2022).
Language models (mostly) know what they know.
\textit{arXiv:2207.05221}.

\bibitem[Kossen et~al.(2024)]{kossen2024semantic}
Kossen, J. et~al. (2024).
Semantic entropy probes.
\textit{arXiv:2406.15927}.

\bibitem[Kuhn et~al.(2023)]{kuhn2023semantic}
Kuhn, L. et~al. (2023).
Semantic uncertainty.
\textit{ICLR 2023}.

\bibitem[Lakshminarayanan et~al.(2017)]{lakshminarayanan2017simple}
Lakshminarayanan, B. et~al. (2017).
Simple and scalable predictive uncertainty estimation using deep ensembles.
\textit{NeurIPS 2017}.

\bibitem[Li et~al.(2024)]{li2024atomic}
Li, A. et~al. (2024).
Atomic self-consistency for better long form generations.
\textit{EMNLP 2024}.

\bibitem[Li et~al.(2026)]{li2026reasc}
Li, R. et~al. (2026).
Reliability-aware adaptive self-consistency.
\textit{arXiv:2601.02970}.

\bibitem[Liang et~al.(2025)]{liang2025cisc}
Liang, J. et~al. (2025).
Confidence-weighted self-consistency.
\textit{ACL Findings 2025}.

\bibitem[Lightman et~al.(2023)]{lightman2023lets}
Lightman, H. et~al. (2023).
Let's verify step by step.
\textit{arXiv:2305.20050}.

\bibitem[Luan et~al.(2018)]{luan2018scierc}
Luan, Y. et~al. (2018).
Multi-task identification of entities, relations, and coreference for scientific knowledge graph construction.
\textit{EMNLP 2018}.

\bibitem[Manakul et~al.(2023)]{manakul2023selfcheckgpt}
Manakul, P. et~al. (2023).
{SelfCheckGPT}: Zero-resource black-box hallucination detection.
\textit{EMNLP 2023}.

\bibitem[N\'{e}dellec \& Xia(2026)]{nedellec2026consistency}
N\'{e}dellec, C. and Xia, M. (2026).
Consistency--accuracy correlation in hard-prompted {LLMs} for entity and relation extraction.
\textit{Genomics Informatics}.

\bibitem[Sang \& De~Meulder(2003)]{sang2003conll}
Sang, E.T.K. and De~Meulder, F. (2003).
Introduction to the {CoNLL}-2003 shared task.
\textit{CoNLL 2003}.

\bibitem[Shelmanov et~al.(2021)]{shelmanov2021certain}
Shelmanov, A. et~al. (2021).
How certain is your {T}ransformer?
\textit{EACL 2021}.

\bibitem[Snell et~al.(2024)]{snell2024scaling}
Snell, C. et~al. (2024).
Scaling {LLM} test-time compute.
\textit{arXiv:2408.03314}.

\bibitem[Vashurin et~al.(2025)]{vashurin2025cocoa}
Vashurin, R. et~al. (2025).
{CoCoA}: Combining confidence and consistency for uncertainty quantification.
\textit{NeurIPS 2025}.

\bibitem[Wang et~al.(2023)]{wang2023selfconsistency}
Wang, X. et~al. (2023).
Self-consistency improves chain of thought reasoning in language models.
\textit{ICLR 2023}.

\bibitem[Whitehead et~al.(2025)]{whitehead2025know}
Whitehead, S. et~al. (2025).
Know your limits: A survey on abstention in large language models.
\textit{TACL 2025}.

\bibitem[Xiao \& Wang(2019)]{xiao2019quantifying}
Xiao, Y. and Wang, W. (2019).
Quantifying uncertainties in natural language processing tasks.
\textit{AAAI 2019}.

\bibitem[Ye et~al.(2025)]{ye2025uq}
Ye, J. et~al. (2025).
Uncertainty quantification for {LLMs}: A survey.
\textit{ACL Findings 2025}.

\bibitem[Zhao et~al.(2026)]{zhao2026selfcertainty}
Zhao, Y. et~al. (2026).
Self-certainty best-of-{N}.
\textit{ICLR 2026}.

\end{thebibliography}
}

\end{document}
